\section{Results}

\subsection*{A dispersion-matched null distinguishes genuine topological
structure from point-cloud spread}

Persistent homology summarises a single-cell state space by the
persistence of its one-dimensional loops (max $H_1$), but this statistic
is sensitive to the overall spread of the point cloud: a featureless but
highly dispersed Gaussian cloud can yield large $H_1$ persistence with no
underlying structure. To separate genuine non-Gaussian organisation from
mere dispersion, we compared every observed max $H_1$ against a
covariance-matched Gaussian null --- a synthetic point cloud drawn with
the same mean and covariance as the data but containing no real
topological structure (Methods). For each cohort we report the observed
max $H_1$ with a 90\% bootstrap confidence interval, the null max $H_1$,
their ratio, and the fraction of bootstrap subsamples in which the
observed value exceeds the matched null. This dispersion-matched null is
the central methodological element of this study: a ratio near or below
$1$ indicates that the apparent loop is consistent with dispersion alone,
whereas a ratio above $1$ that is stable across subsamples indicates
genuine excess structure.

\subsection*{Osimertinib exposure raises max $H_1$ in PC9 cells, robustly
across subsamples}

We analysed a PC9 osimertinib time-series with Watermelon lineage tracing
(D0/D3/D7/D14; Methods). Observed max $H_1$ increased with drug exposure,
from $1.41$ (90\% CI $[1.31, 1.60]$) at D0 to $1.82$ $[1.82, 2.15]$ at D3,
$2.67$ $[2.45, 3.04]$ at D7, and $3.78$ $[2.99, 3.78]$ at D14
(\cref{fig:validation}). Because the point estimates carry appreciable
sampling variability, we did not rely on them directly but instead tested
the \emph{direction} of change across bootstrap subsamples. The increase
was directionally robust: D7 exceeded D0, and D14 exceeded D0, in
$100\%$ of bootstrap subsample pairs. Thus drug exposure raises max $H_1$
in the osimertinib cell-line model reproducibly, independent of the
instability of any single point estimate.

\subsection*{The dispersion-matched null shows the rise reflects genuine
non-Gaussian structure under drug, but not at baseline}

The increase in max $H_1$ is informative only if it exceeds what
dispersion alone would produce. Against the covariance-matched Gaussian
null, the two early conditions were indistinguishable from dispersion:
at D0 the observed-to-null ratio was $0.95$ (null max $H_1 = 1.49$;
observed exceeded the null in only $31\%$ of subsamples), and at D3 the
ratio was $0.91$ (null $2.00$; $18\%$ of subsamples). In both cases the
observed loop is no larger than that of a structureless Gaussian cloud
with the same covariance --- i.e.\ at baseline, max $H_1$ is dispersion.
Under drug, the relationship inverted. At D7 the ratio rose to $1.31$
(null $2.04$), with the observed value exceeding the matched null in
$100\%$ of bootstrap subsamples, and at D14 the ratio was $1.42$
(null $2.66$) with the null exceeded in $98\%$ of subsamples
(\cref{fig:validation}). The emergence of excess, non-Gaussian
topological structure is therefore specific to the drug-exposed
timepoints rather than a property of the cell line, and it appears only
after the dispersion contribution has been controlled.

\subsection*{The drug-induced structure replicates in EGFR-mutant
patient-derived xenografts}

To test whether this controlled signal extends beyond cell culture, we
analysed an EGFR-mutant PDX model under matched untreated and
osimertinib-residual conditions (Methods). In the untreated state the
observed max $H_1$ was $2.79$ (90\% CI $[2.23, 4.03]$), only marginally
above the matched Gaussian null (ratio $1.21$; observed exceeded null in
$86\%$ of subsamples). Under residual disease, max $H_1$ rose to $5.46$
$[3.58, 7.34]$, well above the null (ratio $1.95$; observed exceeded null
in $99\%$ of subsamples). The within-model increase was directionally
robust, with the residual condition exceeding the untreated condition in
$97\%$ of bootstrap subsample pairs (\cref{fig:validation}). The PDX thus
reproduces the cell-line pattern --- weak, dispersion-consistent
structure before treatment and strong, non-Gaussian excess structure
under drug --- in an \emph{in vivo} system with stromal, vascular, and
host context.

\subsection*{Negative control: cells do not traverse the loop}

A persistent loop in max $H_1$ does not by itself imply that cells move
\emph{around} that loop. To test the stronger, traversal-based
interpretation directly, we computed circular coordinates (DREiMac) on
the persistent loop and assessed how cells are distributed around it
(Methods). If cells genuinely traversed the loop, their angular
coordinates would spread around the circle; if they do not, the angles
collapse to a narrow arc. Across all cohorts the circular coordinates
collapsed: the angular distribution had a Rayleigh concentration
$R \geq 0.965$, compared with $R = 0.881$ for a control loop occupied by
$12\%$ of cells. Consistent with this, Watermelon lineage clones were
\emph{tighter} on the loop than randomly chosen cells, rather than spread
around it as transit through a cycle would require. We therefore treat
this as a built-in negative control and explicitly do \emph{not} claim
that cells cycle or traverse the loop. The structure we detect is a
\emph{sparse persistent loop occupied by a minority of cells}, not
population-wide reversible cycling; the topological reorganisation is real
but static in the sense relevant here.

\subsection*{Scope of the claim: erlotinib and patient tumours}

The drug-induction signal is specific, and two further cohorts bound its
scope honestly. First, the PC9 erlotinib time-series did not show excess
topological structure at any timepoint: observed max $H_1$ values were
$1.49$ (D0), $1.77$ (D9), and $1.08$ (D11), with dispersion-matched
ratios of $0.93$, $0.99$, and $0.70$ respectively --- at or below the
Gaussian null throughout. Erlotinib thus produces no excess structure
over dispersion, and we report it not as a positive finding but as
evidence for the specificity of the osimertinib signal; the contrast
between the two TKIs is itself part of the evidence that the controlled
signal is drug- and context-dependent rather than a generic feature of
EGFR-TKI exposure.

Second, in EGFR-mutant patient tumours (Maynard cohort), max $H_1$
exceeded the dispersion-matched null at \emph{every} stage, including
treatment-naive tissue: observed values were $5.49$ (treatment-naive),
$6.04$ (persister), and $5.58$ (progressive disease), with ratios of
$1.52$, $1.69$, and $1.51$ respectively. Because excess structure is
present already in untreated tumours, the patient signal cannot be
attributed to drug induction. We interpret it instead as a reflection of
baseline tumour heterogeneity, which produces non-Gaussian topological
structure independent of treatment. Accordingly, we scope the
drug-induction claim to the controlled osimertinib cell-line and PDX
systems --- where baseline structure is dispersion-consistent and the
excess appears only under drug --- and present the patient data as a
heterogeneity caveat rather than as a drug-response claim.

\subsection*{Summary}

Across controlled models, EGFR-TKI (osimertinib) exposure is accompanied
by genuine non-Gaussian reorganisation of the cell-state space: max $H_1$
rises above a covariance-matched Gaussian dispersion null specifically
under drug (osimertinib D7 and D14; PDX residual disease), robustly
across bootstrap subsamples, whereas baseline conditions are consistent
with dispersion alone. The signal is a sparse persistent loop occupied by
a minority of cells and is not traversed by those cells, as the
circular-coordinate negative control demonstrates. Erlotinib shows no
excess structure, and patient tumours show excess structure already at
baseline, together delimiting the claim to drug-induced topological
reorganisation in the controlled osimertinib cell-line and PDX systems.
The dispersion-matched null that makes these distinctions is a
generalisable, dispersion-controlled topological readout applicable to
any single-cell state space.
